\section{Conclusion}
\label{sec:conclusion}

We developed and evaluated a complete 500-species recognition pipeline on a
deterministic iNaturalist subset.  The rewritten primary traditional model
combines colour, texture, HOG, and spatial SIFT-BoVW in a 3,704-D LinearSVC and
reaches \PrimaryTraditionalTestTopOnePercent\% held-out top-1.  ResNet18
trained from scratch reaches 29.88\%, while ImageNet transfer raises the best
ResNet18 result to 67.60\%.  Pretrained EfficientNet-B0 is the strongest model
at \BestTestTopOnePercent\% top-1, \BestTestTopFivePercent\% top-5, and
\BestTestMacroFOnePercent\% macro F1.  A matched ablation finds a 1.38-point
validation gain from stronger augmentation.

The advanced studies qualify that success.  The corruption analysis uses a
distinct 2,424-D, 300-tree handcrafted Random Forest without BoVW as the
traditional comparator and derives every aggregate only after all 64 cells are
regenerated.  Confusion pairs concentrate among visually similar species, and
Grad-CAM examples suggest both meaningful organism localisation and possible
colour/texture shortcuts.  The main lesson is therefore two-sided: transfer
learning and efficient convolutional architecture provide a large practical
advance, but reliability requires evaluation beyond the clean test set.  The
current selection manifest, persisted predictions, and reproducible analyses
provide a sound basis for future higher-resolution, corruption-aware, and
taxonomy-informed methods, while the retrospective Stage~2 locking limitation
remains part of the evidence record.
